\subsection{Robustness and Scope}
\label{sec:results:scope}

Three checks bear on how far the result on the complete series reaches: the removal of same-day expiries, the years 2008 to 2019 before weekly and daily expirations, and the post-publication period. Table~\ref{tab:scope} gives the rates of the rule and of gpt-6-luna on the complete series (\ident{new_full}) and without same-day expiries (\ident{new_ex0dte}) for each period; Table~S10 of the Supplementary Materials gives them by year, and Figure~S6 of the Supplementary Materials shows 2008 to 2019.

\begin{table}[H]
\centering
\caption{Scope checks, SPY: detection rate (percent) of the classifier and of gpt-6-luna on \ident{new_full} (full) and \ident{new_ex0dte} (ex-0DTE), over all windows whose end lies in the period, and the exact McNemar p-value of the two series on the non-overlapping subsample (Section~\ref{sec:methods:statistics}). The post-publication period is the windows that end from 2026-01-02 to 2026-10-09 (194 windows, 165 of them entirely in 2026; Section~\ref{sec:methods:windows}). Rates by year are in Table~S10 of the Supplementary Materials.}
\label{tab:scope}
\footnotesize
\begin{tabular}{@{}lrrrrrr@{}}
\toprule
 & \multicolumn{2}{c}{Classifier} & \multicolumn{2}{c}{gpt-6-luna} & \multicolumn{2}{c}{McNemar p} \\
\cmidrule(lr){2-3}\cmidrule(lr){4-5}\cmidrule(lr){6-7}
Period & full & ex-0DTE & full & ex-0DTE & classifier & gpt-6-luna \\
\midrule
% src: stats.yaml contrasts series_contrast SPY new_full vs new_ex0dte period=2008-2019 (rate_a_pct, rate_b_pct, nonoverlap_30d.p_mcnemar_exact), classifier and llm (arms full_spy, ex0dte_spy)
2008--2019 & 10.2 & 6.5 & 10.2 & 6.4 & 0.375 & 0.375 \\
% src: stats.yaml contrasts series_contrast SPY new_full vs new_ex0dte period=2020-2023 (rate_a_pct, rate_b_pct, nonoverlap_30d.p_mcnemar_exact), classifier and llm (arms full_spy, ex0dte_spy)
2020--2023 & 22.2 & 18.1 & 22.7 & 18.4 & 1.000 & 0.688 \\
% src: stats.yaml contrasts series_contrast SPY new_full vs new_ex0dte period=2024-2025 (rate_a_pct, rate_b_pct, nonoverlap_30d.p_mcnemar_exact), classifier and llm (arms full_spy, ex0dte_spy)
2024--2025 & 12.0 & 14.7 & 12.5 & 15.3 & 1.000 & 1.000 \\
% src: ex0dte_decomposition.csv symbol=SPY scope=year year=2026 (rate_full_pct, rate_ex0dte_pct, llm_rate_full_pct, llm_rate_ex0dte_pct); stats.yaml contrasts series_contrast new_full vs new_ex0dte period=2026 nonoverlap_30d.p_mcnemar_exact; post_publication_2026_summary.csv definition=ends_in_2026 (11.3, 21.1)
2026 (post-publication) & 11.3 & 21.1 & 11.3 & 21.1 & 1.000 & 1.000 \\
\bottomrule
\end{tabular}
\end{table}

\subsubsection{Same-Day Expiries}
\label{sec:results:ablation}
% src: ex0dte_decomposition.csv (symbol SPY, QQQ, IWM; scope year; abs_0dte_share_pct, sign_flip_by_0dte_pct, rate_full_pct, rate_ex0dte_pct, n_full_only, n_ex0dte_only, llm_rate_*); stats.yaml contrasts series_contrast new_full versus new_ex0dte (paired by window end) for periods 2008-2019, 2020-2023, 2024, 2025, 2024-2025, 2026
The share of same-day expiries in the absolute net GEX of SPY grows from 0.9\% in 2008 and 11.5\% in 2013 to 22.1\% in 2019, 25.7\% in 2020 and 45.9\% in 2024, and it stands at 43.9\% in 2025 and 40.8\% in 2026 (the share is the sum of the absolute same-day-expiry contribution over the sum of the absolute net GEX of \ident{new_full} over the days of the year). For QQQ it is 45.5\% in 2025 and for IWM 23.0\%. The sign of the daily net GEX changes when same-day expiries are removed on 11.1\% of the days of 2024 and 13.2\% of the days of 2025 (4.4\% of all days). The effect on the verdict is small and has no constant direction. In 2024 the rule detects 10.3\% of the windows of the complete series and 13.1\% of those without same-day expiries; 9 windows are detected only on the complete series and 16 only without same-day expiries (difference $-$2.8 percentage points; quarter-block bootstrap interval [$-$10.7, 5.5], $p=0.516$; exact McNemar test on the non-overlapping subsample $p=1.000$). The counts for 2025 are 4 and 11. Pooled over 2024 and 2025 the rates are 12.0\% and 14.7\% for the rule (difference $-$2.8 percentage points, interval [$-$8.7, 2.8], $p=0.361$; McNemar on the subsample of 17 windows $p=1.000$) and 12.5\% and 15.3\% for gpt-6-luna ($-$2.8 percentage points, $p=0.359$; subsample $p=1.000$). Removing same-day expiries therefore does not restore the high rates of the published inputs: the rate is 14.7\% without them and 12.0\% with them, against 100.0\% on the published inputs. In 2020 to 2023 the complete series detects more than the ex-0DTE series (22.2\% against 18.1\%, 4.1 percentage points; $p=0.137$ by bootstrap, $p=1.000$ by McNemar on the subsample), and in 2026 it detects less (11.3\% against 21.1\%; 0 windows against 19 in the paired comparison; bootstrap $p=0.633$, McNemar on the subsample $p=1.000$). For QQQ and IWM the rule detects no window under either series, so the ablation has no effect on their verdicts (Section~S3 of the Supplementary Materials). An example window, 2024-03-15, is in Figure~S3 of the Supplementary Materials.

\subsubsection{2008 to 2019}
\label{sec:results:deep}
% src: detection_by_year.csv (classifier and llm rows, series new_full, new_ex0dte, old_bs, years 2008-2019); stats.yaml rates period 2008-2019; ex0dte_decomposition.csv (diff_pp, llm_diff_pp, abs_0dte_share_pct); stats.yaml contrasts series_contrast new_full vs new_ex0dte period 2008-2019 (classifier: windows_only_a 120, windows_only_b 10, diff_pp 3.67893, nonoverlap_30d.p_mcnemar_exact 0.375; llm: 124 and 12); detection_by_year.csv classifier rows SPY 2008-2014 (max_mean_abs_gex_bn: largest before 2014 3.906766 new_full, 3.648529 new_ex0dte, 4.026031 old_bs; n_shape_ok summed over 2008-2013: 772 of 1,482 on new_full) and report.yaml diagnostics.floor_years; windows parquet new_full_SPY (clf_negative_days >= 21 in 703 of the 772)
Before 2014 no window of any of the three series is detected, by the rule or by gpt-6-luna; this includes the 476 windows that end in the financial crisis of 2008 and 2009. No window of these years reaches the magnitude criterion: the largest 30-day mean $|\mathrm{GEX}|$ is \$3.91~billion on \ident{new_full}, \$3.65~billion on \ident{new_ex0dte} and \$4.03~billion on the archive. The zero is therefore a floor effect of the fixed dollar threshold, as for QQQ and IWM (Section~S3 of the Supplementary Materials), and not an absence of persistent signs: 772 of the 1{,}482 windows of \ident{new_full} that end in 2008 to 2013 (52.1\%) meet the persistence and stability criteria, 703 of them with a negative sign, and fail on magnitude alone. Detections start in 2014, the first year in which a window mean reaches \$5~billion (9.1\% on \ident{new_full}), and peak in 2018 on both complete series, at 34.8\% on \ident{new_full} (87 of 250 windows) and 24.0\% on \ident{new_ex0dte}, the year of the 2018 volatility episode; the archive peaks in 2017 (38.8\%). Pooled over 2008 to 2019 the rule detects 10.2\% of the windows of \ident{new_full} (304 of 2{,}990; Wilson 9.1--11.3\%, quarter-block bootstrap 4.9--16.7\%), 6.5\% of those of \ident{new_ex0dte} and 9.8\% of those of the archive. The model gives 10.2\%, 6.4\% and 9.8\% on the same three series.

% src: detection_by_year.csv label_source=classifier scope=year series_key=new_full, symbols SPY, QQQ (from 2011) and IWM (n_shape_ok and n summed over 2008-2013, 2014-2019, 2020-2023, 2024-2025 and 2026; detection rate equal to the share in 2023-2026 for SPY); series_comparison.csv pair new_full_vs_new_ex0dte symbol SPY scope=year mean_abs_a_bn (0.881 in 2009, 8.384 in 2026; ratio 9.5; Table S7 (tab:app:series) 0.88 and 8.38). The earlier argument that GEX grows with the square of the price (a 58-fold factor from the mean spot) is removed: spot is not in a tracked result file (review items C002, C046)
The fixed threshold is not constant in real terms: the yearly mean $|\mathrm{GEX}|$ of \ident{new_full} rises from \$0.88~billion in 2009 to \$8.38~billion in 2026 (Table~S7 of the Supplementary Materials), a factor of 9.5, so that a window of 2026 meets the \$5~billion criterion far more easily than a window of 2009. Without the magnitude criterion, the share of windows of \ident{new_full} with a persistent sign falls over the whole period: for SPY from 52.1\% in 2008 to 2013 to 41.9\% in 2014 to 2019, 28.8\% in 2020 to 2023 and 12.0\% in 2024 and 2025 (11.3\% in 2026), and in the same direction for QQQ (51.0\% in 2011 to 2013, 13.4\% in 2024 and 2025) and IWM (80.0\% and 35.7\%). From 2023 on the magnitude criterion does not bind for SPY, so the share equals the detection rate and the comparison of Section~\ref{sec:results:replay} is unaffected.

The gap between \ident{new_full} and \ident{new_ex0dte} is a consistency check of the ablation arm. It is zero from 2008 to 2013, where neither series reaches the magnitude criterion, so these years cannot test the ablation, and it is positive in every year from 2014 to 2019: 0.8 percentage points in 2014, 6.3 in 2015, up to 12.7 in 2017 and 6.8 in 2019, while the share rises to 22.1\%. The gap is therefore not proportional to the share. The model reproduces it (up to 12.7 percentage points). Over 2008--2019 the complete series detects 3.7 percentage points more than the ex-0DTE series (120 windows detected only with same-day expiries against 10 only without; quarter-block bootstrap $p\le5.0\times10^{-4}$, the floor of 2{,}000 draws; McNemar on the non-overlapping subsample of 100 windows $p=0.375$). Only the bootstrap separates the two series in the decade before 0DTE; the exact test on the non-overlapping subsample, with 4 against 1 discordant windows, cannot ($p=0.375$), whereas the sign of the difference is reversed, without a significant difference, in 2024 and 2025.


\subsubsection{Post-Publication Period}
\label{sec:results:post}
% src: dates 2026-09-14 (gpt-6-luna creation; README line 216, runs.model_created_at) and 2026-10-12 (Alpha Vantage to Cboe switch; decisions log line 207, README line 126); post_publication_2026_summary.csv (SPY; definitions ends_in_2026 and entirely_2026; series new_full, new_ex0dte, call_only, new_full_scaled); post_publication_2026.csv (clf_regime_type of detected windows, clf_avg_magnitude_bn); stats.yaml rates period 2026
We use the label post-publication rather than out-of-sample, because gpt-6-luna was created on 2026-09-14 and its training cutoff is undocumented. The period has 194 windows with ends from 2026-01-02 to 2026-10-09, of which 165 lie entirely in 2026 (ends from 2026-02-13) and 29 still contain days of 2025. The last window end is 2026-10-09 because equity options moved from Alpha Vantage to Cboe delayed quotes on 2026-10-12 (Section~\ref{sec:methods:lineage}). On the complete series the rule detects 22 of 194 (11.3\%) of the windows that end in 2026; without same-day expiries it detects 41 of 194 (21.1\%) (Table~S11 of the Supplementary Materials also gives the 165 windows entirely in 2026, and Figure~S4 the labels of each window). All detected windows of the complete series are persistent negative-GEX windows, with ends between 2026-03-20 and 2026-04-21. The rate of the complete series (11.3\%) is close to the 12.0\% of 2024 and 2025 on the complete series and far from the 100\% that the published inputs gave for those years. The Wilson interval for the windows ending in 2026 is 7.6--16.6\% (quarter-block bootstrap 0.0--20.3\%, 4 blocks). A call-only series would label 194 of the 194 windows as persistent, with a mean $|\mathrm{GEX}|$ of \$21.74~billion. gpt-6-luna agrees with the rule in 100.0\% of the windows on \ident{new_full} and 100.0\% on \ident{new_ex0dte} (0 and 0 windows of disagreement, respectively). The results for QQQ and IWM in 2026 are in Section~S3 of the Supplementary Materials. The training data of gpt-6-luna may overlap this period, so the agreement of the model with the rule here shows consistency with the rule and says nothing about forecasting.
